\documentclass[../../../include/open-logic-section]{subfiles}

\begin{document}

\olfileid[fr]{sth}{z}{milestone}
\olsection{\texorpdfstring{$\Z^-$}{Z⁻} : un jalon}

Nous reviendrons sur \stagesinf{} dans la section suivante. Toutefois,
avec l'axiome de l'infini, nous avons franchi un cap important. Nous
disposons maintenant de tous les axiomes nécessaires à la théorie
$\Zminus$. Plus précisément :

\begin{defn}
La théorie $\Zminus$ a pour axiomes l'extensionnalité, l'union, la
paire, l'ensemble des parties, l'infini et toutes les instances du
schéma de séparation.
\end{defn}

Ce nom désigne la théorie des ensembles de \emph{Zermelo} (\emph{moins}
quelque chose que nous aborderons plus tard). Zermelo mérite cet hommage,
car il a, pour l'essentiel, formulé cette théorie en
\citeyear{Zermelo1908Untersuchungen}.\footnote{Pour des remarques
intéressantes sur l'histoire et les détails techniques, voir
\citet[annexe~A]{Potter2004}.}

Cette théorie est assez puissante pour nous permettre de faire une
immense quantité de mathématiques. En particulier, vous \emph{devriez}
reparcourir la \olref[sfr][][]{part} et vous convaincre que tout ce que
nous y avons fait de manière naïve pourrait être fait plus formellement
dans~$\Zminus$. (Après vous y être essayé quelque temps, vous pourrez
passer directement à la \olref[sth][z][arbintersections]{sec}.)
Désormais, sans autre précision, nous considérerons que nous
travaillons au moins dans $\Zminus$.

\end{document}
